% ============================================================
%  ANNEXES — TAC-TIC IntelliPM
% ============================================================

\chapter{Glossaire technique}
\label{ann:glossaire}

\begin{description}[leftmargin=4cm, style=nextline]

  \item[API]
  Application Programming Interface — interface logicielle permettant à deux
  applications de communiquer entre elles via des protocoles standardisés (REST, GraphQL).

  \item[AUC ROC]
  Area Under the ROC Curve — métrique d'évaluation d'un modèle de classification
  binaire. Une valeur proche de 1 indique un modèle discriminant.

  \item[Bearer Token]
  Jeton d'authentification stateless inclus dans chaque requête HTTP via l'en-tête
  \texttt{Authorization: Bearer <token>}. Utilisé avec Laravel Sanctum.

  \item[BFS]
  Breadth-First Search (parcours en largeur) — algorithme de graphe utilisé ici
  pour détecter les cycles dans le graphe des dépendances de tâches.

  \item[CI/CD]
  Continuous Integration / Continuous Deployment — pratique DevOps consistant à
  automatiser la compilation, les tests et le déploiement à chaque modification
  du code source.

  \item[CPM]
  Critical Path Method — méthode de planification de projet qui identifie la
  séquence de tâches la plus longue déterminant la durée minimale du projet.

  \item[Docker]
  Plateforme de conteneurisation permettant d'empaqueter une application et ses
  dépendances dans un conteneur isolé et reproductible.

  \item[Eloquent]
  ORM (Object-Relational Mapper) de Laravel permettant d'interagir avec la base
  de données MySQL via des objets PHP plutôt que des requêtes SQL brutes.

  \item[FastAPI]
  Framework Python haute performance pour la construction d'APIs RESTful, basé
  sur les annotations de type Python 3.6+ et le standard ASGI.

  \item[GBM]
  Gradient Boosting Machine — algorithme d'ensemble supervisé qui construit
  séquentiellement des arbres de décision pour minimiser une fonction de perte.

  \item[GitHub Actions]
  Outil d'intégration et de déploiement continus intégré à GitHub, permettant
  de définir des workflows automatisés déclenchés par des événements git.

  \item[InnoDB]
  Moteur de stockage MySQL supportant les transactions ACID, les clés étrangères
  et le niveau d'isolation REPEATABLE READ.

  \item[Joblib]
  Bibliothèque Python permettant la sérialisation (sauvegarde/chargement) efficace
  d'objets Python volumineux, comme les modèles Scikit-learn entraînés.

  \item[Laravel Reverb]
  Serveur WebSocket natif de l'écosystème Laravel, permettant le broadcast
  d'événements en temps réel vers les clients connectés.

  \item[Laravel Sanctum]
  Package Laravel fournissant un système d'authentification léger par tokens
  Bearer pour les SPA et les APIs mobiles.

  \item[LLM]
  Large Language Model — modèle de langage massif entraîné sur de grandes quantités
  de texte pour la génération de texte, la traduction et la compréhension du langage
  naturel (ex. Llama 3.3 70B Instruct).

  \item[MAE]
  Mean Absolute Error — erreur absolue moyenne, métrique de régression mesurant
  l'écart moyen entre les valeurs prédites et réelles.

  \item[MoSCoW]
  Méthode de priorisation des exigences en quatre catégories : Must (obligatoire),
  Should (devrait), Could (pourrait), Won't (ne sera pas). Utilisée pour le
  product backlog.

  \item[MSW]
  Mock Service Worker — bibliothèque JavaScript permettant d'intercepter et
  de simuler des requêtes réseau dans les tests unitaires frontend.

  \item[N+1]
  Problème de performance des ORM où N requêtes SQL supplémentaires sont générées
  pour charger des relations (au lieu d'une seule requête avec JOIN).

  \item[ORM]
  Object-Relational Mapper — couche d'abstraction permettant de manipuler une
  base de données relationnelle via des objets du langage de programmation.

  \item[PestPHP]
  Framework de tests PHP moderne basé sur PHPUnit, avec une syntaxe expressive
  de type BDD (Behavior-Driven Development).

  \item[RBAC]
  Role-Based Access Control — modèle de contrôle d'accès où les permissions
  sont attribuées à des rôles, eux-mêmes assignés aux utilisateurs.

  \item[REST]
  Representational State Transfer — style d'architecture pour les services web
  basé sur les méthodes HTTP (GET, POST, PUT, PATCH, DELETE).

  \item[RF]
  Random Forest — algorithme d'ensemble supervisé qui construit un grand nombre
  d'arbres de décision en parallèle pour améliorer la précision et réduire le
  surapprentissage.

  \item[RMSE]
  Root Mean Square Error — racine carrée de l'erreur quadratique moyenne, métrique
  de régression pénalisant davantage les grandes erreurs.

  \item[ROC]
  Receiver Operating Characteristic — courbe représentant le taux de vrais positifs
  en fonction du taux de faux positifs pour différents seuils de classification.

  \item[Sanctum]
  Voir \textit{Laravel Sanctum}.

  \item[Scikit-learn]
  Bibliothèque Python open-source pour le Machine Learning, incluant des algorithmes
  de classification, régression, clustering et prétraitement des données.

  \item[SonarCloud]
  Service cloud d'analyse statique de code source qui mesure la qualité (bugs,
  vulnérabilités, code smells, duplication, couverture de tests).

  \item[SPA]
  Single Page Application — application web dont le rendu s'effectue côté client
  (navigateur) via JavaScript, sans rechargement de page.

  \item[Spatie]
  Entreprise belge créatrice de nombreux packages Laravel open-source, notamment
  \texttt{laravel-permission} utilisé pour le RBAC système.

  \item[Sprint]
  Itération de développement de durée fixe (ici 2 semaines) dans le cadre de
  la méthodologie Scrum, à l'issue de laquelle un incrément fonctionnel est livré.

  \item[Story Points]
  Unité d'estimation relative de l'effort nécessaire pour réaliser une user story,
  basée sur la complexité, l'effort et l'incertitude.

  \item[TanStack Query]
  Bibliothèque React (anciennement React Query) pour la gestion de l'état serveur :
  mise en cache, invalidation, refetching automatique.

  \item[TailwindCSS]
  Framework CSS utilitaire permettant de styliser des composants directement via
  des classes CSS prédéfinies dans le HTML/JSX.

  \item[User Story]
  Description fonctionnelle d'une fonctionnalité du point de vue de l'utilisateur,
  selon le format : \textit{En tant que [rôle], je veux [action] afin de [bénéfice]}.

  \item[Uvicorn]
  Serveur ASGI (Asynchronous Server Gateway Interface) ultra-performant pour Python,
  utilisé pour servir les applications FastAPI en production.

  \item[Vite]
  Outil de build frontend nouvelle génération basé sur les modules ES natifs,
  offrant un démarrage à froid quasi-instantané et un rechargement à chaud (HMR).

  \item[WebSocket]
  Protocole de communication full-duplex permettant des échanges bidirectionnels
  en temps réel entre le navigateur et le serveur sur une connexion TCP persistante.

\end{description}

% ============================================================
\chapter{Product Backlog complet}
\label{ann:product_backlog}

Le tableau complet du Product Backlog — incluant toutes les 34 user stories,
leur priorisation MoSCoW, leurs story points et leur sprint d'affectation —
est présenté dans le Chapitre~II, section~\ref{tab:product_backlog},
page~\pageref{tab:product_backlog}.

% ============================================================
\chapter{Sprint Backlogs détaillés}
\label{ann:sprint_backlogs}

Les sprint backlogs détaillés de chaque sprint sont présentés dans les chapitres
correspondants :

\begin{description}[leftmargin=3cm]
  \item[Sprint 1] Chapitre III — Tableau~\ref{tab:backlog_s1}
  \item[Sprint 2] Chapitre III — Tableau~\ref{tab:backlog_s2}
  \item[Sprint 3] Chapitre III — Tableau~\ref{tab:backlog_s3}
  \item[Sprint 4] Chapitre IV — Tableau~\ref{tab:backlog_s4}
  \item[Sprint 5] Chapitre IV — Tableau~\ref{tab:backlog_s5}
  \item[Sprint 6] Chapitre V — Tableau~\ref{tab:backlog_s6}
\end{description}

% ============================================================
\chapter{Résultats détaillés de l'analyse SonarCloud}
\label{ann:sonarcloud}

Les résultats complets de l'analyse SonarCloud (rapport PDF d'analyse statique,
métriques de couverture par fichier, liste des code smells corrigés) sont
disponibles dans le dépôt GitHub du projet sous le répertoire
\texttt{docs/sonarcloud-report/}.

Le Quality Gate final est synthétisé dans le Chapitre~VI,
section~\ref{tab:sonar_results}.

% ============================================================
\chapter{Résultats de validation des modèles ML}
\label{ann:ml_validation}

Les métriques détaillées de validation des modèles Machine Learning (matrices de
confusion, courbes ROC par fold, importance des 22 variables, courbes
d'apprentissage) sont disponibles dans le Jupyter Notebook
\texttt{ml-service/notebooks/model\_evaluation.ipynb} du dépôt GitHub.

Les indicateurs clés sont synthétisés dans le Chapitre~IV,
tableau~\ref{tab:ml_metrics}.
